\documentclass[10pt,a4paper]{amsart}
\usepackage[margin=19mm]{geometry}
\usepackage{amsmath,amssymb,mathtools}
\usepackage[scr=rsfs,cal=euler]{mathalfa}
\usepackage{xfrac}
\usepackage[unicode,hidelinks,bookmarks=false]{hyperref}
\newcommand{\ie}{i.e.\ }
\newcommand{\cf}{cf.\ }
\newcommand{\resp}{respectively\ }
\newcommand{\Subsection}{\subsection}
\providecommand{\nobreakdash}{\nobreak\hbox{-}}
\setlength{\emergencystretch}{2em}
\multlinegap0pt
\usepackage{amssymb}
\usepackage[shortlabels]{enumitem}
\usepackage{bm}
\usepackage{xcolor}
\newtheorem{lem}{Lemma}
\title{New directions in the study of prime ideals in rational, nilpotent Iwasawa algebras}
\author{Adam Jones, William Woods}
\date{Source: arXiv:2603.26626v1 (CC BY 4.0)}
\begin{document}
\maketitle

\noindent\emph{Source and licence.} New directions in the study of prime ideals in rational, nilpotent Iwasawa algebras. Version arXiv:2603.26626v1 (CC BY 4.0), distributed by arXiv under the Creative Commons Attribution 4.0 International licence, http://creativecommons.org/licenses/by/4.0/, as recorded in the arXiv licence metadata for this exact version and shown on the abstract page https://arxiv.org/abs/2603.26626v1. Full licence notice: \texttt{licenses/arxiv-2603.26626-CC-BY-4.0.txt}.\par\smallskip

\noindent\textit{Editorial scope.} Counted unit: the article's lem C (source0.tex lines 902--904). The article's complete original proof of it is delivered with it (source0.tex lines 906--925, one \texttt{proof} environment of 20 lines). No mathematics is written, repaired or re-derived: the statement and the proof are the article's own lines, in the article's own notation, and no retained line is substituted or re-flowed.  Retained uncounted as the source-local context the proof needs: lines 621--624.  Nothing was omitted from any retained span, so the package carries no omission record.  No retained line contains a citation, so the wrapper carries no bibliography and no bibliography entry.  No retained line contains a figure environment, an image-inclusion command or drawing code, so no image is delivered and the package declares no asset dependency.  Editorial formatting only, in addition: an AMS \texttt{amsart} preamble that restates the article's own declarations and macros used by the retained lines, namely the environments \texttt{lem} on separate counters, together with the standard packages those lines need.  The retained environments print in this document as Lemma 1 in order of appearance, whereas the article prints the same items with the numbering of its own full document.  The complete native source of the article as distributed in the arXiv e-print for this exact version is bundled unchanged as \texttt{sources/197296/source0.tex}.
\begin{equation}\label{eqn: 2star}
\tag{$**$}
\text{for all } x\in \mathfrak{g}:\, \text{if } [x,[x,\mathfrak{a}]] = 0 \text{ then } [x,\mathfrak{a}] = 0.
\end{equation}

\begin{lem}\label{lem: type C}
Suppose $\mathfrak{s}$ is a simple Lie algebra of type $C_n$ for $n\geq 2$, and $\mathfrak{g}=\mathfrak{n}^+$. Then $B(\mathfrak{g})$ is a maximal abelian ideal in $\mathfrak{g}$ of dimension $\binom{n+1}{2}$.
\end{lem}

\begin{proof}

We can take the root system of type $C_n$ to be the the set $\Phi$ of all integer vectors in $\mathbb{R}^n$ of length $\sqrt{2}$, together with all vectors of the form $2\underline{v}$, where $\underline{v}$ is an integer vector of length 1. Take $$\Phi^+=\{e_i-e_j,2e_k,e_\ell+e_m:i,j,n,k,\ell,m\leq n, i<j,\ell\neq m\}$$

So clearly the set $A=\{e_i+e_j:i,j\leq n\}=\{2e_k,e_\ell+e_m:k,\ell,m\leq n,\ell\neq m\}\subseteq\Phi^+$, and $A$ has size $\binom{n}{2}+n=\frac{n(n-1)}{2}+n=\frac{n(n+1)}{2}=\binom{n+1}{2}$.

Moreover, for all $\alpha,\beta\in A$, $\alpha+\beta\notin\Phi$, so $[\mathfrak{s}_\alpha,\mathfrak{s}_\beta]=0$. In fact, if $\alpha\in\Phi^+$ and $\alpha\notin A$ then $\alpha=e_i-e_j$ for some $i<j$, so for any $\beta=e_s+e_r\in A$, $\alpha+\beta$ lies in $\Phi^+$ if and only if $j=s$ or $r$. Hence $\alpha+\beta=e_i+e_s$ or $e_i+e_r$, and thus $\alpha+\beta\in A$. So either $[\mathfrak{s}_\alpha,\mathfrak{s}_\beta]=0$ or $[\mathfrak{s}_\alpha,\mathfrak{s}_\beta]=s_\gamma$ for some $\gamma\in A$. It follows that $\mathfrak{h}=\underset{\alpha\in A}{\bigoplus}\mathfrak{s}_\alpha$ is an abelian ideal of $\mathfrak{g}$.

Now, suppose $\mathfrak{a}$ is an abelian ideal containing $\mathfrak{h}$. Then for any element $x\in\mathfrak{a}$, write $x=\underset{\alpha\in\Phi^+}{\sum}\lambda_\alpha x_\alpha$, where $x_{\alpha}$ is a generator of $\mathfrak{s}_{\alpha}$. Therefore, $$[x,x_{2e_j}]=\underset{i<j\leq n}{\sum}\lambda_{e_i-e_j}\mu_{e_i-e_j,2e_j}x_{e_i+e_j}$$ and since $x_{2e_j}\in\mathfrak{h}\subseteq\mathfrak{a}$ and $\mathfrak{a}$ is abelian, we know that $[x,x_{2e_j}]=0$, and hence $\lambda_{e_i-e_j}=0$ for all $i<j\leq n$. It follows that $x\in\mathfrak{h}$ and thus $\mathfrak{a}=\mathfrak{h}$, so $\mathfrak{h}$ is a maximal abelian ideal of $\mathfrak{g}$. 

So to prove that $\mathfrak{h}=B(\mathfrak{g})$, it remains to show that $\mathfrak{h}$ satisfies \eqref{eqn: 2star}. We will show that for all $x\in\mathfrak{g}$, if $[x,[x,\mathfrak{s}_{2e_k}]]=0$ for all $k\leq n$ then $x\in\mathfrak{h}$, and it will follow immediately that if $[x,[x,\mathfrak{h}]]=0$ then $[x,\mathfrak{h}]=0$ as required.

We may assume that the $x_{\alpha}$-coefficient of $x$ is zero if $\alpha\in A$, so we can write $x=\underset{i<j\leq n}{\sum}\lambda_{i,j}x_{e_i-e_j}$, and we want to prove that $\lambda_{i,j}=0$ for all $i,j$.

For each $k\leq n$, we see $$[x,x_{2e_k}]=\underset{i<j\leq n}{\sum}\lambda_{i,j}[x_{e_i-e_j},x_{2e_k}]=\underset{i< k}{\sum}\lambda_{i,k}\mu_{e_i-e_k,2e_k}x_{e_i+e_k}$$ thus 
\begin{align*}
0=[x,[x,x_{2e_k}]]&=\underset{i< k}{\sum}\lambda_{i,k}\mu_{e_i-e_k,2e_k}[x,x_{e_i+e_k}]=\underset{\ell<m\leq n}{\sum}\underset{i< k}{\sum}\lambda_{\ell,m}\lambda_{i,k}\mu_{e_i-e_k,2e_k}[x_{e_\ell-e_m},x_{e_i+e_k}]\\&=\underset{\ell<i<k}{\sum}\lambda_{\ell,i}\lambda_{i,k}\mu_{e_i-e_k,2e_k}\mu_{e_\ell-e_i,e_i+e_k}x_{e_\ell+e_k}+\underset{\ell,i<k}{\sum}\lambda_{\ell,k}\lambda_{i,k}\mu_{e_i-e_k,2e_k}\mu_{e_\ell-e_k,e_i+e_k}x_{e_\ell+e_i}
\end{align*}

And note that the $\mu$'s in this equation are all non-zero. Suppose $k$ is minimal such that $\lambda_{j,k}\neq 0$ for some $j<k$. Then the first sum in the equation above is zero (since every term is a multiple of $\lambda_{\ell,i}$ for $i<k$), and the $x_{2e_{j}}$ coefficient of the second sum is $\lambda_{j,k}^2\mu_{e_\ell-e_k,2e_k}\mu_{e_{\ell}-e_k,e_\ell+e_k}$. So since this coefficient must be zero we conclude that $\lambda_{j,k}=0$ -- contradiction. Thus we must have that $\lambda_{i,j}=0$ for all $i,j$ as required.\end{proof}

\end{document}
